\documentclass[11pt]{article}

\usepackage{amsmath,amssymb,amsthm}
\usepackage{algorithm}
\usepackage{algorithmic}
\usepackage{fullpage}
\usepackage{bm}

\newtheorem{theorem}{Theorem}
\newtheorem{assumption}{Assumption}
\newtheorem{lemma}{Lemma}
\newtheorem{definition}{Definition}
\newtheorem{corollary}{Corollary}

\newcommand{\R}{\mathbb{R}}
\newcommand{\E}{\mathbb{E}}
\newcommand{\inner}[2]{\left\langle #1,\, #2 \right\rangle}
\newcommand{\norm}[1]{\left\lVert #1 \right\rVert}

\begin{document}

\section*{Algorithm Name}
Deep Frank-Wolfe with Proximal Linearization (DFW-PL)

\section*{Mathematical Formulation}
We consider a composite objective of the form
\[
\Phi(w) \;=\; \E_{z=(x,y)}\big[\ell(h(w;x),y)\big],
\]
where
- $w\in\R^d$ are parameters,
- $h:\R^d\to\R^m$ is a (possibly non-convex) inner network, differentiable with Jacobian $J(w)=\nabla h(w)\in\R^{m\times d}$,
- $\ell:\R^m\times\mathcal{Y}\to\R_+$ is a convex and smooth outer loss in its first argument (the prediction), applied to $u=h(w;x)$.

At iterate $w_t$, define
\[
u_t := h(w_t;x),\qquad J_t := J(w_t),\qquad g_{u,t} := \nabla_u \ell(u_t,y),\qquad g_t := \nabla \Phi(w_t) \;=\; J_t^\top g_{u,t}.
\]
DFW-PL chooses the descent direction $p_t := -g_t$ and an adaptive closed-form step size $\alpha_t$ by minimizing a tight quadratic upper bound on the exact outer loss applied to the \emph{linearized inner model}:
\[
w_{t+1} \;=\; w_t + \alpha_t p_t,\qquad 
\alpha_t \;=\; \frac{\norm{g_t}^2}{L_\ell \norm{J_t p_t}^2 + \mu \norm{p_t}^2}.
\]
Here $L_\ell$ is the smoothness constant of $\ell$ in its first argument, and $\mu\ge 1$ is a curvature buffer arising from the nonlinearity of the inner model (specified precisely in the convergence analysis). The Jacobian-vector product $J_t p_t$ is computed efficiently by automatic differentiation without forming $J_t$ explicitly.

Hyperparameters and constants:
- $L_\ell>0$: Lipschitz gradient constant of $\ell(\cdot,y)$ in its prediction argument.
- $G_\ell>0$: bound on $\|\nabla_u \ell(u,y)\|$ over the trajectory.
- $\beta>0$: Lipschitz constant of the Jacobian of $h$, i.e., $\|J(w)-J(w')\|\le \beta \|w-w'\|$.
- $\sigma_J>0$: bound on $\|J(w)\|$ over the trajectory.
- $\mu:=\max\!\Big\{1,\, G_\ell \beta + \frac{L_\ell \beta^2}{4}\,(\sigma_J G_\ell)^2\Big\}$ (used in the proofs; any larger $\mu$ is valid).

Remarks:
- The outer loss is kept exact; only the inner model is linearized to derive a provable local quadratic upper bound that yields the closed-form $\alpha_t$ above.
- No hand-designed learning-rate schedule is required.

\section*{Pseudocode}
\begin{algorithm}[H]
\caption{Deep Frank-Wolfe with Proximal Linearization (DFW-PL)}
\begin{algorithmic}[1]
\STATE Input: $w_0\in\R^d$; constants $L_\ell>0$, $G_\ell>0$, $\beta>0$, $\sigma_J>0$
\STATE Set $\mu \leftarrow \max\!\Big\{1,\, G_\ell \beta + \frac{L_\ell \beta^2}{4}\,(\sigma_J G_\ell)^2\Big\}$
\FOR{$t=0,1,2,\dots,T-1$}
  \STATE Compute prediction $u_t = h(w_t;x)$ and $g_{u,t} = \nabla_u \ell(u_t,y)$
  \STATE Backprop composite gradient $g_t = J_t^\top g_{u,t}$ (via reverse-mode AD)
  \STATE Set descent direction $p_t = -g_t$
  \STATE Compute Jacobian-vector product $s_t = J_t p_t$ (via forward-mode or reverse-over-reverse)
  \STATE Set step size $\displaystyle \alpha_t = \frac{\|g_t\|^2}{L_\ell \|s_t\|^2 + \mu \|p_t\|^2}$  // closed-form from quadratic upper bound
  \STATE Update $w_{t+1} = w_t + \alpha_t p_t$
  \IF{stopping criterion satisfied (e.g., $\|g_t\| \le \varepsilon$)}
    \STATE break
  \ENDIF
\ENDFOR
\STATE Output: $w_T$
\end{algorithmic}
\end{algorithm}

\section*{Dry Run Example}
We follow the requested template on the scalar function $f(w)=(w-3)^2$. Here the inner model is the identity $h(w)=w$, the outer loss is $\ell(u)=(u-3)^2$, and $L_\ell=2$ (since $\nabla \ell(u)=2(u-3)$ is $2$-Lipschitz). The requested illustration uses a fixed step $\eta=0.1$ for three iterations.

Initialization: $w_0=0$, gradient $g_0 = \nabla f(w_0)=2(w_0-3) = -6$, objective $f(w_0)=9$.

Iteration 1:
\[
\begin{aligned}
g_0 &= 2(0-3) = -6,\\
w_1 &= w_0 - \eta g_0 = 0 - 0.1(-6) = 0.6,\\
f(w_1) &= (0.6-3)^2 = 2.4^2 = 5.76.
\end{aligned}
\]

Iteration 2:
\[
\begin{aligned}
g_1 &= 2(0.6-3) = -4.8,\\
w_2 &= w_1 - \eta g_1 = 0.6 - 0.1(-4.8) = 1.08,\\
f(w_2) &= (1.08-3)^2 = 1.92^2 = 3.6864.
\end{aligned}
\]

Iteration 3:
\[
\begin{aligned}
g_2 &= 2(1.08-3) = -3.84,\\
w_3 &= w_2 - \eta g_2 = 1.08 - 0.1(-3.84) = 1.464,\\
f(w_3) &= (1.464-3)^2 = 1.536^2 = 2.359296.
\end{aligned}
\]

Note: For this example, the DFW-PL closed-form step specializes to $\alpha_t=\frac{1}{2+\mu}$; choosing $\mu=8$ recovers $\alpha_t=0.1$.

\section*{Assumptions}
We work with the following standard conditions.

\begin{assumption}[Outer loss convex and smooth]\label{asmp:ell}
For each $y$, $\ell(\cdot,y)$ is convex and $L_\ell$-smooth in its first argument:
\[
\ell(v,y) \le \ell(u,y) + \inner{\nabla \ell(u,y)}{v-u} + \frac{L_\ell}{2}\|v-u\|^2,\quad \forall u,v\in\R^m.
\]
Moreover, $\|\nabla \ell(u,y)\| \le G_\ell$ for all $u$ encountered along the iterates.
\end{assumption}

\begin{assumption}[Inner model smoothness]\label{asmp:h}
The inner model $h$ is differentiable with $\beta$-Lipschitz Jacobian:
\[
\|J(w)-J(w')\| \le \beta \|w-w'\|,\quad \forall w,w'\in\R^d.
\]
Consequently, for any $w$ and any direction $p$, the second-order Taylor remainder obeys
\[
\|h(w+\alpha p)-h(w)-\alpha J(w)p\| \le \frac{\beta}{2}\alpha^2\|p\|^2,\quad \forall \alpha\ge 0.
\]
\end{assumption}

\begin{assumption}[Bounded Jacobian]\label{asmp:J}
Along the iterates $\{w_t\}$ we have $\|J(w_t)\|\le \sigma_J$ for some $\sigma_J>0$.
\end{assumption}

\begin{assumption}[Lower bounded objective]\label{asmp:lower}
$\inf_{w}\Phi(w) > -\infty$.
\end{assumption}

\section*{Main Convergence Theorem}
We state a deterministic descent and stationarity-rate guarantee.

\begin{theorem}[Monotone descent and stationarity rate]\label{thm:rate}
Under Assumptions \ref{asmp:ell}--\ref{asmp:lower}, consider DFW-PL with $p_t=-g_t$ and
\[
\alpha_t \;=\; \frac{\|g_t\|^2}{L_\ell \|J_t p_t\|^2 + \mu \|p_t\|^2},\qquad \mu := \max\!\Big\{1,\, G_\ell \beta + \frac{L_\ell \beta^2}{4}\,(\sigma_J G_\ell)^2\Big\}.
\]
Then the iterates satisfy the one-step decrease
\[
\Phi(w_{t+1}) \;\le\; \Phi(w_t) \;-\; \frac{1}{2}\cdot \frac{\|g_t\|^4}{L_\ell \|J_t g_t\|^2 + \mu \|g_t\|^2}
\;\le\; \Phi(w_t) \;-\; \frac{1}{2(L_\ell \sigma_J^2+\mu)}\,\|g_t\|^2.
\]
Summing from $t=0$ to $T-1$ yields
\[
\frac{1}{T}\sum_{t=0}^{T-1}\|g_t\|^2 \;\le\; \frac{2\,(L_\ell \sigma_J^2+\mu)}{T}\,\big(\Phi(w_0)-\inf_w \Phi(w)\big).
\]
In particular, $\min_{0\le t\le T-1}\|g_t\|^2 = O\big(1/T\big)$, i.e., the method converges to a first-order stationary point at the standard nonconvex rate with explicit constants.
\end{theorem}

\section*{Theoretical Properties}
- Stability: The step size satisfies $\alpha_t \le 1/\mu \le 1$, ensuring conservative steps when curvature is large.
- Hyperparameter sensitivity: Only $\mu$ aggregates problem constants $(L_\ell,G_\ell,\beta,\sigma_J)$; any larger $\mu$ preserves all guarantees (more conservative).
- Comparison to baselines:
  - Versus SGD with hand-tuned schedules, DFW-PL computes $\alpha_t$ in closed form from curvature proxies $L_\ell\|J_t p_t\|^2$ and $\mu\|p_t\|^2$.
  - Versus Frank-Wolfe on parameter domains, DFW-PL performs a Frank-Wolfe-like linearization in the \emph{model space} while keeping the outer loss exact.
  - Versus first-order methods with global $L$, DFW-PL adapts locally through $J_t p_t$ without backtracking.

\section*{Discussion}
- When $h$ is linear ($\beta=0$), we may set $\mu=1$; the step size simplifies and the bound matches classical smooth convex descent guarantees.
- For common losses (squared, logistic), the same derivation applies with the same $L_\ell$.

\section*{Preliminaries (Definitions and Lemmas)}
We first collect two standard lemmas.

\begin{lemma}[Smoothness of $\ell$]\label{lem:ell_smooth}
Under Assumption \ref{asmp:ell}, for any $a,r\in\R^m$,
\[
\ell(a+r,y) \;\le\; \ell(a,y) + \inner{\nabla \ell(a,y)}{r} + \frac{L_\ell}{2}\|r\|^2.
\]
\end{lemma}

\begin{lemma}[Second-order remainder bound for $h$]\label{lem:h_remainder}
Under Assumption \ref{asmp:h}, for any $w\in\R^d$, direction $p\in\R^d$, and $\alpha\ge 0$,
\[
R(\alpha;w,p) \;:=\; h(w+\alpha p)-h(w)-\alpha J(w)p,\qquad \|R(\alpha;w,p)\| \;\le\; \frac{\beta}{2}\alpha^2\|p\|^2.
\]
\end{lemma}

\noindent We also record a bound following from Assumptions \ref{asmp:ell} and \ref{asmp:J}:
\begin{equation}\label{eq:gw_bound}
\|g_t\| = \|J_t^\top g_{u,t}\| \le \|J_t\|\,\|g_{u,t}\| \le \sigma_J G_\ell.
\end{equation}

\section*{Main Proof (Step-by-Step Derivation)}
We prove Theorem \ref{thm:rate}. Fix an iteration $t$, write $w=w_t$, $J=J_t$, $u=h(w)$, $g_u=\nabla\ell(u,y)$, $g=\nabla\Phi(w)=J^\top g_u$, and consider any direction $p\in\R^d$ and scalar $\alpha\ge 0$.

Define the actual next prediction and its linear approximation:
\[
u^+ \;:=\; h(w+\alpha p),\qquad \tilde{u} \;:=\; u + \alpha J p,\qquad R := u^+ - \tilde{u}.
\]
By Lemma \ref{lem:h_remainder}, $\|R\| \le \frac{\beta}{2}\alpha^2\|p\|^2$. We bound the one-step change in $\Phi$.

\noindent\textbf{Part I: Keep the outer loss exact while using the linearized inner model.}
\begin{align}
\Phi(w+\alpha p) &= \E\big[\ell(u^+,y)\big] \nonumber\\
&= \E\big[\ell(\tilde{u}+R,y)\big]. \nonumber
\end{align}
Apply Lemma \ref{lem:ell_smooth} at $a=\tilde{u}$ and $r=R$:
\begin{align}
\ell(\tilde{u}+R,y) 
&\le \ell(\tilde{u},y) + \inner{\nabla\ell(\tilde{u},y)}{R} + \frac{L_\ell}{2}\|R\|^2. \tag*{[Step 1]}
\end{align}
Take expectation and use $\|\nabla\ell(\tilde{u},y)\|\le G_\ell$:
\begin{align}
\Phi(w+\alpha p)
&\le \E[\ell(\tilde{u},y)] + \E\big[\inner{\nabla\ell(\tilde{u},y)}{R}\big] + \frac{L_\ell}{2}\E\|R\|^2 \nonumber\\
&\le \E[\ell(\tilde{u},y)] + G_\ell\,\E\|R\| + \frac{L_\ell}{2}\E\|R\|^2. \tag*{[Step 2]: Cauchy-Schwarz and gradient bound}
\end{align}
Using Lemma \ref{lem:h_remainder} and dropping the expectation (deterministic upper bound):
\begin{align}
\E\|R\| \le \frac{\beta}{2}\alpha^2\|p\|^2,\qquad
\E\|R\|^2 \le \left(\frac{\beta}{2}\alpha^2\|p\|^2\right)^2 = \frac{\beta^2}{4}\alpha^4\|p\|^4. \tag*{[Step 3]}
\end{align}
Thus
\begin{align}
\Phi(w+\alpha p) \le \E[\ell(\tilde{u},y)] + \frac{G_\ell \beta}{2}\alpha^2\|p\|^2 + \frac{L_\ell \beta^2}{8}\alpha^4\|p\|^4. \tag*{[Step 4]}
\end{align}

\noindent\textbf{Part II: Upper bound $\E[\ell(\tilde{u},y)]$ via smoothness of $\ell$.}
Apply Lemma \ref{lem:ell_smooth} at $a=u$ and $r=\alpha J p$:
\begin{align}
\ell(\tilde{u},y) = \ell(u+\alpha J p,y)
&\le \ell(u,y) + \alpha\,\inner{\nabla\ell(u,y)}{J p} + \frac{L_\ell}{2}\alpha^2 \|J p\|^2. \tag*{[Step 5]}
\end{align}
Taking expectation and recognizing $\E[\ell(u,y)]=\Phi(w)$ and $\E[\nabla\ell(u,y)]=g_u$:
\begin{align}
\E[\ell(\tilde{u},y)] 
&\le \Phi(w) + \alpha\,\inner{g_u}{J p} + \frac{L_\ell}{2}\alpha^2 \|J p\|^2. \tag*{[Step 6]}
\end{align}
Use $\inner{g_u}{J p} = \inner{J^\top g_u}{p} = \inner{g}{p}$:
\begin{align}
\E[\ell(\tilde{u},y)] 
&\le \Phi(w) + \alpha\,\inner{g}{p} + \frac{L_\ell}{2}\alpha^2 \|J p\|^2. \tag*{[Step 7]}
\end{align}

\noindent\textbf{Part III: Combine Parts I and II, and dominate the $\alpha^4$ term.}
From [Step 4] and [Step 7],
\begin{align}
\Phi(w+\alpha p) 
&\le \Phi(w) + \alpha\,\inner{g}{p} + \frac{L_\ell}{2}\alpha^2 \|J p\|^2 + \frac{G_\ell \beta}{2}\alpha^2\|p\|^2 + \frac{L_\ell \beta^2}{8}\alpha^4\|p\|^4. \tag*{[Step 8]}
\end{align}
By \eqref{eq:gw_bound}, for the DFW-PL direction we will set $p=-g$ and then $\|p\|=\|g\|\le \sigma_J G_\ell$; thus
\begin{align}
\|p\|^4 \le (\sigma_J G_\ell)^2 \|p\|^2. \tag*{[Step 9]}
\end{align}
Moreover, we preselect any $\mu$ such that
\begin{align}
\mu \;\ge\; G_\ell \beta + \frac{L_\ell \beta^2}{4}(\sigma_J G_\ell)^2 \qquad\Longrightarrow\qquad
\frac{G_\ell \beta}{2}\alpha^2\|p\|^2 + \frac{L_\ell \beta^2}{8}\alpha^4\|p\|^4 \;\le\; \frac{\mu}{2}\alpha^2\|p\|^2, \quad \forall \alpha\in[0,1]. \tag*{[Step 10]}
\end{align}
Combining [Step 8]--[Step 10], for any $\alpha\in[0,1]$,
\begin{align}
\Phi(w+\alpha p) 
&\le \Phi(w) + \alpha\,\inner{g}{p} + \frac{1}{2}\alpha^2\Big(L_\ell \|J p\|^2 + \mu \|p\|^2\Big). \tag*{[Step 11]}
\end{align}

\noindent\textbf{Part IV: Choose $p=-g$ and minimize the quadratic upper bound in $\alpha$.}
Set $p^\star=-g$ so that $\inner{g}{p^\star} = -\|g\|^2$ and $Jp^\star = -Jg$. Then the right-hand side of [Step 11] is a convex quadratic in $\alpha\in[0,1]$:
\begin{align}
Q(\alpha) := \Phi(w) - \alpha \|g\|^2 + \frac{1}{2}\alpha^2\Big(L_\ell \|J g\|^2 + \mu \|g\|^2\Big). \tag*{[Step 12]}
\end{align}
Its unconstrained minimizer over $\alpha\ge 0$ is
\begin{align}
\alpha_{\mathrm{quad}} \;=\; \frac{\|g\|^2}{L_\ell \|J g\|^2 + \mu \|g\|^2}. \tag*{[Step 13]: derivative zero condition}
\end{align}
Note that $\alpha_{\mathrm{quad}}\le \frac{1}{\mu}\le 1$ by $\mu\ge 1$, so $\alpha_{\mathrm{quad}}\in[0,1]$. Plugging $\alpha_{\mathrm{quad}}$ into $Q$ yields 
\begin{align}
Q(\alpha_{\mathrm{quad}}) 
&= \Phi(w) - \frac{1}{2}\cdot \frac{\|g\|^4}{L_\ell \|J g\|^2 + \mu \|g\|^2}. \tag*{[Step 14]: standard quadratic completion}
\end{align}
Therefore, with $p=-g$ and $\alpha=\alpha_{\mathrm{quad}}$, [Step 11] implies
\begin{align}
\Phi(w+\alpha_{\mathrm{quad}} p) \le Q(\alpha_{\mathrm{quad}}) 
= \Phi(w) - \frac{1}{2}\cdot \frac{\|g\|^4}{L_\ell \|J g\|^2 + \mu \|g\|^2}. \tag*{[Step 15]}
\end{align}
Finally, using $\|J g\|\le \|J\|\,\|g\| \le \sigma_J \|g\|$,
\begin{align}
\frac{\|g\|^4}{L_\ell \|J g\|^2 + \mu \|g\|^2} 
\;\ge\; \frac{\|g\|^4}{(L_\ell \sigma_J^2 + \mu)\|g\|^2}
\;=\; \frac{\|g\|^2}{L_\ell \sigma_J^2 + \mu}. \tag*{[Step 16]}
\end{align}
Combining [Step 15] and [Step 16] proves the one-step decrease:
\begin{align}
\Phi(w_{t+1}) \le \Phi(w_t) - \frac{1}{2}\cdot \frac{\|g_t\|^4}{L_\ell \|J_t g_t\|^2 + \mu \|g_t\|^2}
\le \Phi(w_t) - \frac{1}{2(L_\ell \sigma_J^2 + \mu)}\,\|g_t\|^2. \tag*{[Step 17]}
\end{align}
Summing [Step 17] over $t=0,\dots,T-1$ and using $\Phi(w_T)\ge \inf_w \Phi(w)$ yields
\begin{align}
\sum_{t=0}^{T-1}\|g_t\|^2 \;\le\; 2\,(L_\ell \sigma_J^2 + \mu)\,\big(\Phi(w_0)-\inf_w \Phi(w)\big). \tag*{[Step 18]}
\end{align}
Dividing by $T$ completes the proof of Theorem \ref{thm:rate}. \qed

\section*{Rate Derivation (With Constants)}
From [Step 18],
\[
\frac{1}{T}\sum_{t=0}^{T-1}\|\nabla \Phi(w_t)\|^2 
\;\le\; \frac{2\,(L_\ell \sigma_J^2 + \mu)}{T}\,\big(\Phi(w_0)-\Phi^\star\big),
\]
where $\Phi^\star:=\inf_w \Phi(w)$. With the explicit choice
\[
\mu \;=\; \max\!\Big\{1,\, G_\ell \beta + \frac{L_\ell \beta^2}{4}\,(\sigma_J G_\ell)^2\Big\},
\]
the constants are fully specified. In particular, if $h$ is linear ($\beta=0$), we may set $\mu=1$ and obtain
\[
\frac{1}{T}\sum_{t=0}^{T-1}\|\nabla \Phi(w_t)\|^2 
\;\le\; \frac{2\,(L_\ell \sigma_J^2 + 1)}{T}\,\big(\Phi(w_0)-\Phi^\star\big).
\]

\section*{Special Cases}
- Linear inner model ($\beta=0$): $h(w)=Aw+b$, $J\equiv A$, no remainder term; the step size reduces to
\[
\alpha_t = \frac{\|g_t\|^2}{L_\ell \|A g_t\|^2 + \|g_t\|^2} \le 1,
\]
and the descent bound is
\[
\Phi(w_{t+1}) \le \Phi(w_t) - \frac{1}{2(L_\ell \|A\|^2+1)}\|g_t\|^2.
\]
- Identity inner model with squared loss: $h(w)=w$, $\ell(u)=(u-y)^2$, $L_\ell=2$, $\beta=0$, $\mu=1$. Then
\[
\alpha_t = \frac{\|g_t\|^2}{2\|g_t\|^2 + \|g_t\|^2} = \frac{1}{3},\quad
w_{t+1} = w_t - \tfrac{1}{3} g_t,
\]
and $\Phi(w_{t+1})\le \Phi(w_t) - \frac{1}{2(2+1)}\|g_t\|^2$.

\end{document}